% Revize výukové přípravy, výchozí soubor GitHub Work-MAT, 91. hodina.
% Zdroje: Realisticky.cz 020210–020213; rozšiřující úlohy vlastní
% Úlohy jsou vlastní / úlohy výchozí přípravy, případně nově didakticky obměněné.
\documentclass[11pt,a4paper]{article}
\input{../../sablony/pripravy/priprava.sty}
\usepackage{tikz}
\nastavPripravu{1. ročník}{Rovnice a nerovnice}{91}{Slovní úlohy na lineární rovnice}
\begin{document}
\blok{Sestavení rovnice ze situace}
Nestačí dosazovat do stále stejného vzorce „počáteční stav + sazba $\cdot$ čas“. Dnes rozlišujeme \textbf{jednu měnící se veličinu, porovnání dvou nabídek a setkání dvou pohybujících se objektů}. Vše lze vyřešit jednou lineární rovnicí.
\blok{Jak postupovat}
\begin{enumerate}[label=\arabic*)]
\item Označ neznámou a její jednotku; urči přípustné hodnoty.
\item Z každé věty sestav výraz popisující tutéž veličinu na obou stranách rovnosti.
\item Vyřeš rovnici a vrať kořen k původnímu významu; proveď věcnou zkoušku.
\end{enumerate}
\blok{Společně řešené příklady}
\begin{prikladReseniBox}[title={Příklad 1}]
Taxi účtuje nástupní cenu 50 Kč a 18 Kč za kilometr. Kolik kilometrů ujelo při ceně 230 Kč?
\tcblower
\[\begin{aligned}
50+18k&=230 \\
18k&=180 \\
k&=10\,\mathrm{km}
\end{aligned}\]
Při $k=10$ vychází $50+18\cdot10=230$ Kč.
\end{prikladReseniBox}
\begin{prikladReseniBox}[title={Příklad 2}]
Dva telefonní tarify: A stojí 120 Kč měsíčně a 4 Kč/min, B stojí 300 Kč a 2 Kč/min. Při kolika minutách stojí stejně? Který je levnější před tímto okamžikem?
\tcblower
\[\begin{aligned}
120+4m&=300+2m \\
2m&=180 \\
m&=90\,\mathrm{min}
\end{aligned}\]
Zkouška: oba stojí 480 Kč. Při $m=0$ je levnější A, proto A je levnější pro $m<90$, B pro $m>90$.
\end{prikladReseniBox}
\begin{prikladReseniBox}[title={Příklad 3}]
Dva cyklisté vyjíždějí současně proti sobě z míst vzdálených 63 km. Jejich rychlosti jsou 18 a 24 km/h. Za jak dlouho se setkají?
\tcblower
\[\begin{aligned}
18t+24t&=63 \\
42t&=63 \\
t&=1{,}5\,\mathrm{h}
\end{aligned}\]
Zkouška: $27+36=63$ km. Každý příspěvek představuje ujetou vzdálenost.
\end{prikladReseniBox}
\blok{Samostatná práce}
\begin{zadaniReseniBox}[title={Samostatná práce 1}]
Pronájem nářadí stojí 200 Kč plus 45 Kč za hodinu. Za kolik hodin zaplatíme 650 Kč?
\tcblower
\[\begin{aligned}
200+45h&=650 \\
45h&=450 \\
h&=10\,\mathrm h
\end{aligned}\]
Zkouška: $200+45\cdot10=650$ Kč.
\end{zadaniReseniBox}
\begin{zadaniReseniBox}[title={Samostatná práce 2}]
Dvě firmy vyrábějí stejnou součástku. První má hotovo 80 kusů a vyrábí 12 kusů/h, druhá má 140 kusů a vyrábí 8 kusů/h. Za kolik hodin budou mít stejně kusů?
\tcblower
\[\begin{aligned}
80+12t&=140+8t \\
4t&=60 \\
t&=15\,\mathrm h
\end{aligned}\]
Obě mají po 15 h $260$ kusů.
\end{zadaniReseniBox}
\begin{zadaniReseniBox}[title={Samostatná práce 3}]
Obvod obdélníku je 54 cm, délka je o 5 cm větší než šířka. Určete oba rozměry.
\tcblower
\[\begin{aligned}
2x+2(x+5)&=54 \\
4x+10&=54 \\
x&=11\,\mathrm{cm} \\
x+5&=16\,\mathrm{cm}
\end{aligned}\]
Zkouška: $2\cdot(11+16)=54$ cm.
\end{zadaniReseniBox}
\blok{Volitelná rezerva}
\begin{zadaniReseniBox}[title={Rozšíření 1}]
Vlak vyráží v 8:00 rychlostí 70 km/h. Auto vyráží po stejné trase v 9:00 rychlostí 105 km/h. Kdy dojede vlak?
\tcblower
\[\begin{aligned}
105t&=70(t+1) \\
35t&=70 \\
t&=2\,\mathrm h
\end{aligned}\]
Čas $t$ měříme od 9:00. Auto vlak dojede v 11:00.
\end{zadaniReseniBox}
\blok{Domácí úkol}
\begin{zadaniReseniBox}[title={Domácí úkol}]
\begin{enumerate}[label=\alph*),itemsep=2pt]
\item V nádrži je 12 l kapaliny, která ubývá rychlostí $0{,}7\,\mathrm{l/h}$. Za jak dlouho zbude 5 l?
\item Pronájem stojí 350 Kč plus 25 Kč za hodinu. Kdy je účet 900 Kč?
\item Zisk při výrobě $q$ kusů je $100q-1500$ Kč. Kolik kusů je nutné pro zisk 2500 Kč?
\item Tarif A stojí $150+4m$ Kč, tarif B $250+2m$ Kč. Při jakém počtu minut $m$ jsou ceny stejné?
\item Vozidlo A jede rychlostí 45 km/h; vozidlo B vyjede po stejné trase o 30 minut později rychlostí 60 km/h. Za kolik hodin od startu B dojede vozidlo A?
\item Obvod obdélníku je 66 cm. Délka je o 3 cm větší než šířka. Určete rozměry.
\end{enumerate}
\tcblower
\begin{enumerate}[label=\alph*),itemsep=1pt]
\item
\[\begin{aligned}
12-0{,}7t&=5 \\
0{,}7t&=7 \\
t&=10\,\mathrm h
\end{aligned}\]\par Zkouška: za 10 hodin ubyde 7 l.
\item
\[\begin{aligned}
350+25t&=900 \\
25t&=550 \\
t&=22\,\mathrm h
\end{aligned}\]
\item
\[\begin{aligned}
100q-1500&=2500 \\
100q&=4000 \\
q&=40\ \text{kusů}
\end{aligned}\]
\item
\[\begin{aligned}
150+4m&=250+2m \\
2m&=100 \\
m&=50\,\mathrm{min}
\end{aligned}\]
\item
\[\begin{aligned}
60t&=45(t+0{,}5) \\
15t&=22{,}5 \\
t&=1{,}5\,\mathrm h
\end{aligned}\]\par Od startu A uplynou 2 hodiny.
\item
\[\begin{aligned}
2x+2(x+3)&=66 \\
4x+6&=66 \\
x&=15\,\mathrm{cm} \\
x+3&=18\,\mathrm{cm}
\end{aligned}\]\par Obvod $2(15+18)=66$ cm.
\end{enumerate}
\end{zadaniReseniBox}
\poznamka{45 minut: úvodní model 6, tři společné typy 21, samostatné procvičování 14, shrnutí 4. Rezerva je nad rámec; vyžaduj jednotky a zkoušku v příběhu.}
\end{document}
